\chapter{مقدمه}
\label{ch:intro}

بخش بزرگی از دانش هر سازمانی جایی ذخیره نشده که بشود از آن پرس‌وجو گرفت و در متن آزاد پراکنده است. برای پاسخ به پرسشی که در نگاه اول ساده می‌آید، مثلاً اینکه مجری فلان پروژه چه شرکتی بوده یا فلان شخص تابع کدام کشور است، باید ده‌ها سند را دستی خواند و کنار هم گذاشت.

گراف دانش راه‌حل شناخته‌شده این وضعیت است: متن به گره و یال تبدیل می‌شود و پرسش، به‌جای خواندن، با یک پرس‌وجو پاسخ می‌گیرد. مسئله این است که ساختن خود گراف از متن، کار ساده‌ای نیست.

\section{بیان مسئله}
\label{sec:problem}

آنچه در این پژوهش دنبال شد، استخراج رابطه در سطح سند\LTRfootnote{Document-level Relation Extraction} است. ورودی یک متن چندجمله‌ای است و خروجی مجموعه‌ای از سه‌تایی‌ها که هرکدام به جمله منبع خود گره خورده‌اند. دو چیز کار را سخت می‌کند.

اولی اینکه رابطه معمولاً در یک جمله جمع نشده و مدل باید چند جمله را دنبال کند و از زنجیره هم‌مرجعی رد شود. دومی که به نظر ما جدی‌تر است، این است که مدل زبانی رابطه نادرست می‌سازد. برای یک گراف دانش این خطا از خطای معمول بدتر است، چون یال غلط یک بار وارد می‌شود و بعد برای همیشه در پایگاه‌داده می‌ماند و به پرسش‌های بعدی هم سرایت می‌کند.

\section{مشارکت‌های این کار}
\label{sec:contributions}

ایده محوری این است که به مدل زبانی اعتماد نکنیم و خروجی‌اش را از چند صافی رد کنیم. پنج جزء زیر روی هم این کار را انجام می‌دهند.

\subsection{هستان‌شناسی مقیّد}
\label{sec:c-ontology}

سه‌تایی تنها وقتی پذیرفته می‌شود که ترکیب «نوع فاعل، رابطه، نوع مفعول» آن در فهرست مجاز باشد و هر چیز دیگری بی‌چون‌وچرا کنار گذاشته می‌شود. این صافی در اجرا روی \lr{Re-DocRED} یک‌پنجم خروجی مدل، یعنی $20.3\%$، را رد کرد و همان‌طور که در بخش \ref{sec:supervision} نشان می‌دهیم، بهایی که بابتش پرداختیم تنها $0.21\%$ از بازخوانی بود.

\subsection{شاهد اجباری}
\label{sec:c-evidence}

هر یال باید عین جمله منبع را با خود حمل کند و یال بی‌شاهد اصلاً وارد گراف نمی‌شود. حاصلش این است که هیچ ادعایی در گراف نمی‌ماند که نشود آن را تا متن اصلی دنبال کرد.

\subsection{لایه بستار قاعده‌محور}
\label{sec:c-closure}

مدل زبانی آنچه را می‌گوید که در جمله نوشته شده، اما گراف دانش باید آنچه را هم نگه دارد که از خود گراف نتیجه می‌شود. سه قاعده روی گراف اجرا می‌شوند و حقیقت‌های تازه می‌سازند. نکته طراحی اینجاست که هر حقیقت استنتاجی شاهد مقدمه‌های خود را به ارث می‌برد، پس قانون بند قبل زیر پا گذاشته نمی‌شود.

\subsection{خط لوله تکرارپذیر}
\label{sec:c-idempotent}

سند تازه اضافه کنید و دوباره اجرا کنید؛ نه سند تکراری ساخته می‌شود، نه قطعه، نه سه‌تایی، نه گره و نه یال. این را با کلید پایدار \lr{entity\_key} برای گره‌ها، \lr{fact\_id} مبتنی بر \lr{sha1} برای یال‌ها و کش در سطح \lr{chunk\_sha} به دست آوردیم.

\subsection{ورودی چندقالبی}
\label{sec:c-input}

علاوه بر \lr{PDF} و \lr{DOCX} و \lr{TXT}، تصویر هم پذیرفته می‌شود و با مدل بینایی رونویسی می‌گردد.

\section{ساختار گزارش}
\label{sec:structure}

فصل \ref{ch:related} کارهای پیشین و مبانی لازم را مرور می‌کند. فصل \ref{ch:method} معماری سامانه، هستان‌شناسی، لایه بستار و مهندسی لازم برای مقیاس صدهزار سندی را شرح می‌دهد. فصل \ref{ch:dataset} داده و آرایش آزمایش و معیارها را معرفی می‌کند و توضیح می‌دهد سامانه دقیقاً از چه نظارتی استفاده می‌کند. فصل \ref{ch:results} نتایج سه مرحله کار و تحلیل خطا را می‌آورد. فصل \ref{ch:casestudy} همان سامانه را روی اسناد فارسی اجرا می‌کند و فصل \ref{ch:conclusion} محدودیت‌ها و جمع‌بندی را می‌گوید.
